\begin{frame}{Pipeline di esecuzione}
  \begin{center}
    \begin{adjustbox}{max width=\textwidth, max height=0.85\textheight, keepaspectratio}
      \begin{tikzpicture}[
          node distance=1.8ex and 3em,
          base/.style={draw, thick, rounded corners=3pt, inner sep=6pt, align=center, font=\small, text width=13.5em},
          % Colonna sinistra ora "rossina"
          phase-startup/.style={base, fill=red!5, draw=red!50},
          % Colonna destra "verdina"
          phase-exec/.style={base, fill=green!5, draw=green!55},
          term/.style={draw=gray!50, thick, fill=gray!10, rounded corners=4pt, minimum height=3ex, inner sep=5pt, align=center, font=\footnotesize\ttfamily},
          lbl/.style={font=\footnotesize\bfseries},
          arr/.style={-latex, thick, gray!70!black, rounded corners=4pt}
        ]

        % Nodo iniziale
        \node[term] (start) {apiguard run};

        % COLONNA SINISTRA (Bloccanti - Rossa)
        \node[phase-startup, below=5ex of start, xshift=-9em] (p1) {\textbf{Fase~1:} Inizializzazione\\[0.5ex] \textcolor{red!70!black}{\scriptsize\ttfamily ConfigurationError}};
        \node[phase-startup, below=of p1] (p2) {\textbf{Fase~2:} OpenAPI Discovery\\[0.5ex] \textcolor{red!70!black}{\scriptsize\ttfamily OpenAPILoadError}};
        \node[phase-startup, below=of p2] (p3) {\textbf{Fase~3:} Costruzione Contesti};
        \node[phase-startup, below=of p3] (p4) {\textbf{Fase~4:} Test Discovery e DAG\\[0.5ex] \textcolor{red!70!black}{\scriptsize\ttfamily DAGCycleError}};

        % COLONNA DESTRA (Non bloccanti - Verde)
        \node[phase-exec, below=5ex of start, xshift=9em] (p5) {\textbf{Fase~5:} Esecuzione\\[0.5ex] \textcolor{red!70!black}{\scriptsize\ttfamily TestResult(ERROR), isolato}};
        \node[phase-exec, below=of p5] (p6) {\textbf{Fase~6:} Teardown LIFO\\[0.5ex] \textcolor{red!70!black}{\scriptsize\ttfamily TeardownError, WARNING}};
        \node[phase-exec, below=of p6] (p7) {\textbf{Fase~7:} Report};

        % Etichette ancorate con i colori abbinati
        \node[lbl, text=red!70!black, anchor=south west, yshift=1ex] at (p1.north west) {Fasi 1-4 (bloccanti)};
        \node[lbl, text=green!70!black, anchor=south west, yshift=1ex] at (p5.north west) {Fasi 5-7 (non bloccanti)};

        % Nodo finale: spostato esattamente sotto la Fase 7
        \node[term, below=3ex of p7] (end) {exit code: 0 / 1 / 2 / 10};

        % FLUSSO FRECCE

        % 1. Discesa e ingresso pulito nella Fase 1
        \draw[arr] (start.south) -- ++(0, -2ex) -| (p1.north);

        \draw[arr] (p1) -- (p2);
        \draw[arr] (p2) -- (p3);
        \draw[arr] (p3) -- (p4);

        % 2. Salto interno (risale il canale vuoto)
        \draw[arr] (p4.east) -- ++(1.5em,0) |- (p5.west);

        \draw[arr] (p5) -- (p6);
        \draw[arr] (p6) -- (p7);

        % 3. Uscita dritta verticale dalla Fase 7 all'exit code
        \draw[arr] (p7.south) -- (end.north);

      \end{tikzpicture}
    \end{adjustbox}
  \end{center}
\end{frame}